%======================================================================
\section{How to compare TTNs and DNNs}\label{sec:compare}
%======================================================================

The comparison should not be tied to independent input blocks.  Independence is a useful
\emph{diagonalization assumption}: it makes products of left and right cut features orthogonal,
turns the population loss into a tensor-product norm, and gives a genuine Schmidt decomposition
at every physical cut.  None of the basic empirical questions requires it.  Loss trajectories,
function-space projections, representation probes, Jacobian observables and causal interventions
can all be defined under an arbitrary joint input law.  What changes is the geometry in which the
teacher's cut features have to be compared.

This section therefore separates three regimes.
\begin{enumerate}[label=(\roman*),nosep]
\item A \emph{physical TTN with product data}: different leaves receive independent input blocks.
      This is the clean special case in which the data-weighted HSVD is a functional HSVD.
\item A \emph{physical TTN with dependent data}: different leaves receive different coordinate
      blocks, but the blocks have an arbitrary joint distribution.  Cut features remain functions
      of the two sides, but their products are correlated and need a joint Gram correction.
\item A \emph{diagonal, or computational, TTN}: every formal leaf receives the same input.  The
      joint law of the leaves is supported on the diagonal.  This includes the tied TTN that is
      exactly equivalent to a deep monomial or polynomial network.  Its copy-cut HSVD describes a
      chosen computation, not a canonical decomposition of the input--output function.
\end{enumerate}

For each regime we distinguish three levels of evidence.  At the \emph{output} level, does the
student predictor contain the teacher computation?  At the \emph{representation} level, can the
corresponding feature be decoded from an internal state?  At the \emph{use} level, is the downstream
network sensitive to that code and causally reliant on it?  These levels should not be collapsed:
a feature can be present but unused, or used in a distributed form that is not linearly localized.

%----------------------------------------------------------------------
\subsection{Teacher--student setting and input laws}\label{sec:compare-setting}
%----------------------------------------------------------------------

\paragraph{Physical teacher.}
Fix a binary tree with leaves $a\in[N]$.  Leaf $a$ receives a block $x^{(a)}\in\mathcal X_a$
through a fixed embedding $\xi_a:\mathcal X_a\to\R^{d_a}$.  The teacher tensor $T^\star$ defines
\begin{equation}\label{eq:teacher}
 f^\star(x)=\Big\langle T^\star,\ \bigotimes_{a=1}^N\xi_a\big(x^{(a)}\big)\Big\rangle .
\end{equation}
The input is now allowed to have an arbitrary joint law
\begin{equation}\label{eq:joint-blocks}
 X=(X^{(1)},\ldots,X^{(N)})\sim\nu_x,
 \qquad
 \E\Big\|\bigotimes_{a=1}^N\xi_a(X^{(a)})\Big\|^2<\infty .
\end{equation}
The product-law case is the additional assumption
\begin{equation}\label{eq:independent-blocks}
 \nu_x=\mu_1\otimes\cdots\otimes\mu_N .
\end{equation}
The blocks need not be Gaussian, centred, whitened or identically distributed in either case.

For a set $A\subseteq[N]$, write
\[
 X_A=(X^{(a)})_{a\in A},\qquad
 \Xi_A(X_A):=\bigotimes_{a\in A}\xi_a(X^{(a)}),
\]
and let $\nu_A$ denote the marginal law of $X_A$.  A bond $e$ of the teacher tree determines a
physical cut $(A_e\mid A_e^c)$ and a matricization $M_e^\star=T^{\star\langle A_e\rangle}$, so that
\begin{equation}\label{eq:teacher-cut}
 f^\star(x)=\Xi_{A_e}(x_{A_e})^\top M_e^\star\Xi_{A_e^c}(x_{A_e^c}).
\end{equation}
This pointwise identity does not require independence.

\paragraph{Diagonal teacher.}
For a deep monomial network of exponent $K$ and depth $L$, $N=K^L$ formal leaves receive the
same input and
\begin{equation}\label{eq:diagonal-teacher}
 f^\star(x)=\langle T^\star,x^{\otimes N}\rangle .
\end{equation}
For a general polynomial network, $x$ is replaced by the homogenized input $\bar x=(1,x)$.
The formal leaf law is $(\Delta_N)_\#\mu$, where $\Delta_N(x)=(x,\ldots,x)$; it is maximally
dependent.  The equivalence with a standard monomial network requires the structured cores of
Proposition~\ref{prop:monomial-layer}, tied by height.  A generic height-tied multilinear core is a
more general homogeneous polynomial layer and need not factor through one weight matrix followed
by a coordinatewise monomial.

\paragraph{Students and time.}
Every student is trained on the same examples and target.  A DNN and a global TTN need not use the
same raw learning rate: their parameter metrics and function-space speeds differ.  The optimizer
family, tuning budget and time convention should be fixed before looking at the feature curves.
For a local polynomial proxy, by contrast, the original parameters and optimizer state are shared,
so identical updates are the correct comparison.  For independent global students, report update
count and integrated learning-rate time, and permit at most one pre-registered global calibration
such as matching the initial function-space speed; feature-specific or post hoc time warps would
make the surrogate hypothesis unfalsifiable.

%----------------------------------------------------------------------
\subsection{A cut decomposition under an arbitrary joint law}\label{sec:dw-hsvd}
%----------------------------------------------------------------------

Fix a physical cut $e$ and abbreviate
\[
 L_e:=\Xi_{A_e}(X_{A_e}),\qquad R_e:=\Xi_{A_e^c}(X_{A_e^c}),
\]
with marginal Gram matrices
\[
 \Sigma_{e,L}:=\E[L_eL_e^\top],\qquad
 \Sigma_{e,R}:=\E[R_eR_e^\top].
\]
These are Gram matrices of the full subtree embeddings.  Under dependence they generally do not
factor into Kronecker products of leaf Gram matrices and may contain high-order joint moments of
the raw data.

\begin{definition}[Marginally weighted cut SVD]\label{def:dw-hsvd}
Let
\[
 \widetilde L_e:=(\Sigma_{e,L}^{1/2})^+L_e,
 \qquad
 \widetilde R_e:=(\Sigma_{e,R}^{1/2})^+R_e,
\]
and take the compact SVD
\begin{equation}\label{eq:cut-svd}
 \widetilde M_e^\star
 :=\Sigma_{e,L}^{1/2}M_e^\star\Sigma_{e,R}^{1/2}
 =U_e^\star S_e^\star V_e^{\star\top},
 \qquad
 S_e^\star=\diag(\sigma_{e,1}^\star,\ldots,\sigma_{e,r_e}^\star).
\end{equation}
The associated left and right cut functions are
\begin{equation}\label{eq:singular-functions}
 \phi_{e,q}^\star(x_{A_e})=u_{e,q}^{\star\top}\widetilde L_e(x_{A_e}),
 \qquad
 \psi_{e,q}^\star(x_{A_e^c})=v_{e,q}^{\star\top}\widetilde R_e(x_{A_e^c}).
\end{equation}
We call~\eqref{eq:cut-svd} a \emph{marginally weighted cut SVD}.  Under the product law it is the
data-weighted HSVD used in the previous version of this section.
\end{definition}

\begin{proposition}[What the cut SVD does and does not give]\label{prop:dw-hsvd}
Under~\eqref{eq:joint-blocks}, for every physical cut $e$:
\begin{enumerate}[label=(\roman*),nosep]
\item The families $(\phi_{e,q}^\star)_q$ and $(\psi_{e,q}^\star)_q$ are orthonormal in
      $L^2(\nu_{A_e})$ and $L^2(\nu_{A_e^c})$, respectively.
\item The teacher has the pointwise expansion
\begin{equation}\label{eq:cut-expansion}
 f^\star(x)=\sum_{q=1}^{r_e}\sigma_{e,q}^\star
 \phi_{e,q}^\star(x_{A_e})\psi_{e,q}^\star(x_{A_e^c})
 \qquad \nu_x\text{-a.s.}
\end{equation}
\item If $X_{A_e}$ and $X_{A_e^c}$ are independent, the product functions
      $\phi_{e,i}^\star\psi_{e,j}^\star$ are orthonormal in $L^2(\nu_x)$, and
      \eqref{eq:cut-expansion} is the Schmidt decomposition of $f^\star$ across the cut.
\item Under general dependence, the product functions need not be orthogonal.  The numbers
      $\sigma_{e,q}^\star$ are then coefficient singular values, not the $L^2(\nu_x)$ strengths of
      the products, and truncating~\eqref{eq:cut-expansion} need not be the best rank-$r$
      approximation in population loss.
\item If all leaves are independent, $\Sigma_{e,L}$ and $\Sigma_{e,R}$ factor into products of
      leaf Gram matrices for every bond.  One local whitening of the open legs then produces a
      single weighted TTN whose SVDs are nested across the tree.  Under generic dependence,
      whitening is cut-specific and the cut SVDs need not form one nested HSVD.
\end{enumerate}
\end{proposition}

\begin{proof}
The vectors $L_e$ and $R_e$ lie almost surely in the ranges of their Gram matrices.  Hence
$L_e^\top M_e^\star R_e=\widetilde L_e^\top\widetilde M_e^\star\widetilde R_e$ almost surely.
The nonzero left singular vectors lie in $\operatorname{range}(\Sigma_{e,L}^{1/2})$, so
$\E[\phi_{e,q}^\star\phi_{e,q'}^\star]=\delta_{qq'}$; similarly on the right.  Inserting the SVD
proves~\eqref{eq:cut-expansion}.  Independence factorizes the expectation of a product of a left
and a right function, proving (iii).  The remaining statements follow because, without this
factorization, the joint fourth moments
$\E[\phi_i\phi_k\psi_j\psi_l]$ are not products of marginal second moments, and because a
cut-specific square root $\Sigma_{e,L}^{1/2}$ cannot in general be absorbed into independent leaf
maps.
\end{proof}

\begin{remark}[Computation-relative versus function-relative modes]\label{rem:joint-identifiability}
For a fixed teacher tensor, the cut SVD is gauge invariant.  Under dependent data it need not be an
invariant of the on-support function.  If the products of the left and right embedding coordinates
are linearly dependent on the support of $\nu_x$, two different coefficient matrices can define the
same function almost surely.  This is especially important on the diagonal, where antisymmetric
multilinear extensions vanish.  The modes in~\eqref{eq:cut-svd} should therefore be called planted
or computation-relative targets unless the joint product Gram defined below is nonsingular.
\end{remark}

%----------------------------------------------------------------------
\subsection{Joint-Gram atoms and output observables}\label{sec:atoms}
%----------------------------------------------------------------------

For a cut $e$, collect the marginally orthonormal factors into column vectors
$\Phi_e^\star=(\phi_{e,1}^\star,\ldots,\phi_{e,r_e}^\star)^\top$ and
$\Psi_e^\star=(\psi_{e,1}^\star,\ldots,\psi_{e,r_e}^\star)^\top$.  Define the complete product
feature vector
\begin{equation}\label{eq:product-dictionary}
 Z_e(X):=\Psi_e^\star(X_{A_e^c})\otimes\Phi_e^\star(X_{A_e})\in\R^{r_e^2},
\end{equation}
and let $s_e^\star:=\vect(S_e^\star)$.  Then~\eqref{eq:cut-expansion} is simply
$f^\star(X)=Z_e(X)^\top s_e^\star$.

\begin{definition}[Joint product Gram and Gram-corrected projection]\label{def:coupling}
For any joint input law, define
\begin{equation}\label{eq:joint-gram}
 G_e:=\E[Z_e(X)Z_e(X)^\top],
 \qquad
 b_e[f]:=\E[Z_e(X)f(X)]=\vect\!\left(\E[f(X)\Phi_e^\star\Psi_e^{\star\top}]\right).
\end{equation}
The raw matrix inside the vectorization is still called the \emph{coupling matrix}; it is defined
without independence, but it is not a coordinate matrix unless $G_e=I$.  The orthogonal projection
of $f\in L^2(\nu_x)$ onto
$\mathcal D_e:=\operatorname{span}\{Z_{e,ij}:i,j\le r_e\}$ is
\begin{equation}\label{eq:gram-projection}
 (\Pi_e f)(x):=Z_e(x)^\top\beta_e[f],
 \qquad
 \beta_e[f]:=G_e^+b_e[f].
\end{equation}
\end{definition}

\begin{proposition}[Exact output decomposition for arbitrary joint inputs]\label{prop:error-decomposition}
For every $f\in L^2(\nu_x)$:
\begin{enumerate}[label=(\roman*),nosep]
\item Equation~\eqref{eq:gram-projection} is the $L^2(\nu_x)$ orthogonal projection onto
      $\mathcal D_e$, and
      $\|\Pi_e f\|^2=b_e[f]^\top G_e^+b_e[f]$.
\item If $b_e^\star:=b_e[f^\star]=G_es_e^\star$, then
\begin{equation}\label{eq:joint-error}
 \|f-f^\star\|^2
 =\underbrace{(b_e[f]-b_e^\star)^\top G_e^+(b_e[f]-b_e^\star)}_{\text{error inside the teacher factor-product space}}
 +\underbrace{\|f\|^2-b_e[f]^\top G_e^+b_e[f]}_{\text{outside error}}.
\end{equation}
\item If $G_e$ is nonsingular, $\beta_e[f^\star]=s_e^\star$.  If it is singular, only the
      equivalence class $s_e^\star+\ker G_e$ is observable from data; no output statistic can
      identify its individual coordinates.
\end{enumerate}
\end{proposition}

\begin{proof}
If $c\in\ker G_e$, then $c^\top Z_e=0$ almost surely, so $c^\top b_e[f]=0$ and
$b_e[f]\in\operatorname{range}G_e$.  The normal equations for least squares in the dictionary
$Z_e$ are therefore solved by $G_e^+b_e[f]$.  This proves (i).  Since $f^\star\in\mathcal D_e$,
$f-f^\star=\Pi_e(f-f^\star)+(I-\Pi_e)f$ is an orthogonal decomposition, which gives~\eqref{eq:joint-error}.
The last claim is the standard identifiability statement for a singular Gram matrix.
\end{proof}

The teacher SVD supplies spectral blocks $B$ of equal or nearly equal singular values.  The
basis-invariant target associated with such a block is
\begin{equation}\label{eq:block-atom}
 m_{e,B}^\star(x):=\sum_{q\in B}\sigma_{e,q}^\star
 \phi_{e,q}^\star(x_{A_e})\psi_{e,q}^\star(x_{A_e^c}).
\end{equation}
For a simple singular value, this is the usual atom $m_{e,q}^\star$.  Let
$\mathcal A_e:=\operatorname{span}\{m_{e,B}^\star:B\}$ and let $P_{\mathcal A_e}$ be its
orthogonal projector.  Since $f^\star\in\mathcal A_e\subseteq\mathcal D_e$, there is a second,
basis-invariant decomposition
\begin{equation}\label{eq:three-space-decomposition}
 \|f-f^\star\|^2
 =\underbrace{\|P_{\mathcal A_e}(f-f^\star)\|^2}_{\text{teacher-mode error}}
 +\underbrace{\|(\Pi_e-P_{\mathcal A_e})f\|^2}_{\text{alternative pairings inside }\mathcal D_e}
 +\underbrace{\|(I-\Pi_e)f\|^2}_{\text{outside error}}.
\end{equation}
This replaces the previous entrywise ``within/across/outside'' split, which is nonnegative only
when the product features are orthogonal.

To define acquisition coefficients in correlated data, collect the canonical block atoms into a
vector $m_e^\star(X)$ and let
\begin{equation}\label{eq:block-regression}
 H_e:=\E[m_e^\star m_e^{\star\top}],
 \qquad
 \alpha_e[f]:=H_e^+\E[m_e^\star f].
\end{equation}
If $H_e$ is nonsingular, $\alpha_e[f^\star]=\mathbf 1$ and $\alpha_{e,B}(f_t)$ is the partial
least-squares acquisition coefficient of block $B$, controlling for correlations with the other
teacher blocks.  If $H_e$ is singular or badly conditioned, individual block coefficients are not
scientifically identifiable; report the subspace errors in~\eqref{eq:three-space-decomposition}
and merge unstable blocks.

\begin{corollary}[What independence buys]\label{prop:atoms}
If $X_{A_e}$ and $X_{A_e^c}$ are independent, then $G_e=I$.  Consequently:
\begin{enumerate}[label=(\roman*),nosep]
\item the products $\phi_{e,i}^\star\psi_{e,j}^\star$ are an orthonormal dictionary;
\item $\sigma_{e,q}^{\star2}$ is the $L^2$ energy of the simple atom, and the block atoms are
      mutually orthogonal;
\item $\alpha_{e,q}[f]=\langle f,\phi_{e,q}^\star\psi_{e,q}^\star\rangle/\sigma_{e,q}^\star$,
      recovering the acquisition coefficient used in~\eqref{eq:acquisition};
\item truncating the Schmidt expansion gives the best rank-$r$ approximation across the cut by
      the Eckart--Young theorem;
\item the old Frobenius coupling-matrix decomposition is the specialization of
      \eqref{eq:joint-error} and~\eqref{eq:three-space-decomposition}.
\end{enumerate}
\end{corollary}

%----------------------------------------------------------------------
\subsection{Student modes: product-law HSVD and dependent-law alternatives}\label{sec:student-fsvd}
%----------------------------------------------------------------------

Under~\eqref{eq:independent-blocks},
$L^2(\nu_x)=L^2(\nu_{A_e})\otimes L^2(\nu_{A_e^c})$, so every square-integrable student predictor
has a genuine Schmidt decomposition
\begin{equation}\label{eq:student-fsvd}
 f(x)=\sum_{q\ge1}s_{e,q}(f)\,\phi^f_{e,q}(x_{A_e})\psi^f_{e,q}(x_{A_e^c}).
\end{equation}
The singular values and leading subspaces can be compared with the teacher, and estimated with the
crossed-grid matrix
$[f(x_k^{\mathrm L},x_l^{\mathrm R})/\sqrt{n_Ln_R}]_{kl}$ built from independent marginal
samples.  This is a strong, architecture-independent observable in the product-law benchmark.

For dependent inputs there is no natural isometric identification of $L^2(\nu_x)$ with the Hilbert
tensor product of the two marginal spaces.  A crossed grid samples
$\nu_{A_e}\otimes\nu_{A_e^c}$ rather than $\nu_x$ and may evaluate the model on combinations that
never occur.  Consequently, a product-measure SVD is an off-distribution diagnostic, not a
functional HSVD of the training problem.  Three replacements are available:
\begin{enumerate}[label=(\alph*),nosep]
\item the teacher-anchored joint-Gram projections~\eqref{eq:joint-error}, which require only
      ordinary joint samples and are the default here;
\item a finite Galerkin interaction operator
      $B_e[f]=\E[f(X)\Phi^D(X_{A_e})\Psi^D(X_{A_e^c})^\top]$ in pre-specified marginal
      dictionaries; its SVD describes task-weighted cross-dependence, but not an orthogonal
      expansion of $f$;
\item if a trustworthy conditional generator is available, conditional resampling from
      $\nu(dx_{A_e^c}\mid x_{A_e})$ to estimate operators under the actual joint law.
\end{enumerate}
The second and third choices should be labelled explicitly; neither restores the Eckart--Young or
nested-HSVD conclusions without additional assumptions.

%----------------------------------------------------------------------
\subsection{The tied repeated-input case}\label{sec:repeated-input}
%----------------------------------------------------------------------

The diagonal TTN is not a small perturbation of the independent-block setting.  Its formal leaves
are perfectly correlated.  Let $\bar X=(1,X)$ for a polynomial network, and put
$z_N(X)=\bar X^{\otimes N}$.  For two coefficient tensors $T$ and $T'$ define
$f_T(x)=\langle T,z_N(x)\rangle$.  Then
\begin{equation}\label{eq:moment-operator}
 \langle f_T,f_{T'}\rangle_{L^2(\mu)}
 =\vect(T)^\top\mathcal G_N\vect(T'),
 \qquad
 \mathcal G_N:=\E[z_N(X)z_N(X)^\top].
\end{equation}
The operator $\mathcal G_N$ contains moments up to order $2N$.  It is singular on the full tensor
space because $z_N(X)$ is symmetric, and
$f_T=f_{\Sym T}$.  Thus the population loss depends on the symmetrized polynomial coefficient,
whereas the HSVD of the particular unrolled tensor $T$ can change when one changes the off-diagonal
multilinear extension.

This leads to a two-track comparison.
\begin{enumerate}[label=(\roman*),nosep]
\item \emph{Function track.}  Compare the student and teacher in an orthonormal polynomial basis
      for the actual input law.  For Gaussian inputs, Hermite tensors separate the Wiener-chaos
      orders; for a general distribution, use an empirically orthonormalized polynomial dictionary
      up to the required degree.  Coefficients, subspace projections and loss contributions in this
      basis are function-level observables.
\item \emph{Computation track.}  Track the bond singular values and subtree singular functions of
      the fixed, architecture-induced tied TTN.  These quantities are gauge invariant for that
      computation and can be compared with hidden states of the neural network, but they are not
      invariants of the diagonal function.  The subtree of height $\ell$ equals the hidden state
      $h_\ell(x)$ for the exact tied construction, so hidden-state subspaces and the computational
      HSVD are the natural internal observables.
\end{enumerate}
For Gaussian $X$, the expectations in~\eqref{eq:moment-operator} are exactly computable by Wick
contractions.  They do not reduce to products of leaf covariances, and deeper models require higher
moments.  For general data, use Monte Carlo estimates on natural samples.  A crossed-grid estimator
would instead probe an arbitrary off-diagonal extension and should not be used as a function-level
metric.

%----------------------------------------------------------------------
\subsection{Representation observables}\label{sec:representation}
%----------------------------------------------------------------------

Let $h_\ell(x)\in\R^{d_\ell}$ be a student representation.  For any scalar teacher target
$\zeta\in L^2(\nu_x)$ with nonzero variance, define the held-out linear acquisition score
\begin{equation}\label{eq:probe}
 R_h^2[\zeta]
 :=1-\frac{\min_{w,b}\E[(\zeta(X)-w^\top h(X)-b)^2]}{\Var\zeta}
 =\frac{\Cov(h,\zeta)^\top\Cov(h)^+\Cov(h,\zeta)}{\Var\zeta}.
\end{equation}
For a vector target $Z$, use the average squared canonical correlation
\begin{equation}\label{eq:vector-probe}
 R_h^2[Z]
 :=\frac1{\rank\Cov(Z)}\operatorname{tr}
 \big(\Cov(Z)^+\Cov(Z,h)\Cov(h)^+\Cov(h,Z)\big).
\end{equation}
These definitions require no independence.  They are invariant under invertible affine changes of
hidden coordinates and, for vector targets, under changes of basis inside a degenerate teacher
subspace.

For a physical cut, probe for the left factors $\phi_{e,q}^\star$, the right factors
$\psi_{e,q}^\star$, the block atoms $m_{e,B}^\star$, and, when identifiable, the simple products
$\phi_{e,q}^\star\psi_{e,q}^\star$.  For a diagonal TTN, probe instead for teacher hidden-state
subspaces, computational singular functions, and orthogonal polynomial output modes.

\begin{remark}[Dependence creates a predictability confound]\label{rem:probe-meaning}
When the blocks are dependent, a representation that only receives $X_{A_e}$ can predict a function
of $X_{A_e^c}$ through statistical correlation.  A high score for a right factor therefore does not
show that the representation computed information from the right side.  Report the best baseline
from the raw available variables, such as $R^2_{X_{A_e}}[\psi_{e,q}^\star]$, and, when feasible,
probe the conditional residual
$\psi_{e,q}^\star-\E[\psi_{e,q}^\star\mid X_{A_e}]$.  More generally, compare teacher targets with
matched-complexity control targets and with initialization.  Linear decodability establishes
accessibility, not use or causal necessity.
\end{remark}

%----------------------------------------------------------------------
\subsection{Use and causal reliance}\label{sec:use}
%----------------------------------------------------------------------

Write $f_\theta=f_{>\ell}\circ h_\ell$ and whiten the current representation by
$h_\ell^{\mathrm w}=(\Cov(h_\ell)^{1/2})^+(h_\ell-\E h_\ell)$.  For a represented target $\zeta$,
let
\[
 u_\ell[\zeta]
 :=\frac{\E[h_\ell^{\mathrm w}\zeta]}{\|\E[h_\ell^{\mathrm w}\zeta]\|}
\]
be its linear code direction.  The gradient of the output with respect to whitened coordinates is
$\partial_\ell f=\Cov(h_\ell)^{1/2}\nabla f_{>\ell}(h_\ell)$.

\begin{definition}[Downstream sensitivity]\label{def:sensitivity}
Let $\mathrm{AGOP}_\ell=\E[\partial_\ell f\,\partial_\ell f^\top]$.  Define
\[
 \operatorname{sens}_\ell[\zeta]
 :=u_\ell[\zeta]^\top\mathrm{AGOP}_\ell u_\ell[\zeta],
 \qquad
 \overline{\operatorname{sens}}_\ell[\zeta]
 :=\frac{\operatorname{sens}_\ell[\zeta]}
 {\operatorname{tr}(\mathrm{AGOP}_\ell)/\rank\Cov(h_\ell)}.
\]
For a target subspace with orthonormal code matrix $U$, use
$\operatorname{tr}(U^\top\mathrm{AGOP}_\ell U)$.
\end{definition}

Sensitivity is local and unsigned.  A signed gain is obtained by injecting
$h_\ell^{\mathrm w}(x)\mapsto h_\ell^{\mathrm w}(x)+\varepsilon\zeta(x)u$ and differentiating the
chosen output coefficient or Gram-projected block acquisition at $\varepsilon=0$.  Under dependence,
this derivative is still estimated on ordinary joint samples.

\begin{definition}[Erasure and injection]\label{def:interventions}
For a scalar or vector target with code subspace $U$, \emph{linear erasure} replaces
$h_\ell^{\mathrm w}$ by $(I-UU^\top)h_\ell^{\mathrm w}$.  \emph{Injection} adds
$\varepsilon\zeta(x)u$, or the corresponding vector-valued signal along an orthonormal family.
The intervention may be applied once or maintained during subsequent training.
\end{definition}

The erasure removes all affine information about the target at the checkpoint, but nonlinear
information can survive and the network can re-encode the feature later.  Measure both the immediate
change in loss and output projections and the subsequent change in acquisition times.  Use
random-direction, matched-norm and sham controls.  These interventions remain the strongest common
observables across product, dependent and repeated-input regimes.

%----------------------------------------------------------------------
\subsection{What is exact for TTN students?}\label{sec:ttn-students}
%----------------------------------------------------------------------

The answer depends on the input law.
\begin{enumerate}[label=(\roman*),nosep]
\item \emph{Physical TTN, independent leaves.}  Population inner products factor into Kronecker
      products of leaf Gram matrices.  The student's functional HSVD is its data-weighted HSVD,
      coupling matrices and loss decompositions are obtained by deterministic contractions, and
      second-moment messages close on the tree.
\item \emph{Physical TTN, dependent leaves.}  The algebraic cut factorization
      $g_\vartheta=h_e(X_{A_e})^\top n_e(X_{A_e^c})$ still holds, but $h_e$ and $n_e$ are correlated.
      A bond state may predict environment features through input correlations.  Exact population
      quantities require the joint Gram $G_e$ or higher mixed moments and generally do not reduce
      to leafwise second moments.  Monte Carlo evaluation of~\eqref{eq:joint-gram} remains unbiased;
      exact message passing is available only when the data law itself has compatible tractable
      structure.
\item \emph{Tied diagonal TTN.}  Every subtree sees the same input, and the height-$\ell$ subtree
      state is the hidden state of the corresponding monomial or polynomial network.  Shared-core
      gradients are sums over all occurrences.  Function-space loss is governed by the moment
      operator~\eqref{eq:moment-operator}; the computational HSVD can be contracted exactly from
      the cores, but its singular values are not population energies unless this is separately
      established.
\end{enumerate}

%----------------------------------------------------------------------
\subsection{Acquisition times and estimation}\label{sec:compare-protocol}
%----------------------------------------------------------------------

If a block coefficient $\alpha_{e,B}(t)$ in~\eqref{eq:block-regression} is identifiable, define
\begin{equation}\label{eq:acq-times}
 \tau^{\mathrm{out}}_{e,B}(\delta)
 :=\inf\{t:\alpha_{e,B}(s)\ge\delta\text{ for all }s\in[t,t+\Delta]\}.
\end{equation}
If it is not identifiable, define acquisition through a monotone subspace statistic, for example the
fraction of teacher-mode energy captured by $P_{\mathcal A_e}f_t$, and report the conditioning of
$H_e$.  Representation acquisition times are defined from the held-out probe scores in
\eqref{eq:probe}--\eqref{eq:vector-probe}; use several thresholds and treat unobserved crossings as
right-censored.

A practical protocol is:
\begin{enumerate}[label=\arabic*.,nosep,leftmargin=*]
\item Fix the teacher, joint input law and natural audit distribution.  Pre-register the primary cuts,
      blocks and repeated-input polynomial modes.
\item Estimate all Gram matrices on data disjoint from training and testing.  Regularize small
      eigenvalues and report effective rank and condition number; merge modes whose separation is
      below estimation error.
\item On a held-out joint audit set, estimate $b_e[f_t]$, $G_e$, the two terms of
      \eqref{eq:joint-error}, and the three subspace errors in~\eqref{eq:three-space-decomposition}.
      Use sample splitting for quadratic functionals.  Do not recombine marginals unless the product
      law is part of the experiment.
\item Fit representation probes on separate train, validation and test splits.  Under dependence,
      include raw-side and conditional-predictability baselines.  For repeated inputs, include
      polynomial/Hermite coefficient trajectories and teacher hidden-subspace probes.
\item Measure sensitivity, erasure and injection responses at the same checkpoints.  Compare not
      only unperturbed acquisition order but the shifts caused by the same teacher-feature
      intervention in the TTN and DNN.
\item Pair teacher instances, datasets and random seeds.  Report loss and feature trajectories in
      the pre-specified time convention, together with initialization scale, optimizer and
      parameterization; do not hide mismatches behind post hoc time warping.
\end{enumerate}

%----------------------------------------------------------------------
\subsection{Summary: what survives without independence}\label{sec:compare-summary}
%----------------------------------------------------------------------

\begin{table}[t]
\centering
\small
\renewcommand{\arraystretch}{1.2}
\begin{tabular}{p{0.23\textwidth}p{0.31\textwidth}p{0.36\textwidth}}
\hline
\raggedright Object & \raggedright Arbitrary joint inputs & \raggedright Extra conclusion under independent blocks \tabularnewline
\hline
\raggedright Loss and function distance
& \raggedright Defined exactly in $L^2(\nu_x)$; sample estimates are immediate.
& \raggedright Coefficient-space metric factorizes into leaf Gram matrices. \tabularnewline
\raggedright Teacher cut targets
& \raggedright Marginally weighted cut SVD; products require the joint Gram $G_e$ and may be non-identifiable.
& \raggedright One nested data-weighted HSVD; product atoms are orthonormal and canonical in function space. \tabularnewline
\raggedright Output acquisition
& \raggedright Gram-corrected projections, block regressions and subspace errors.
& \raggedright Raw coupling entries are orthogonal coordinates; additive per-atom loss and Eckart--Young truncation. \tabularnewline
\raggedright Student modes
& \raggedright Teacher-anchored projections or stated Galerkin/conditional operators.
& \raggedright Genuine functional Schmidt decomposition and crossed-grid estimator. \tabularnewline
\raggedright Representation and use
& \raggedright Probes, AGOP sensitivity, erasure, injection and acquisition times all generalize; correlation baselines are required.
& \raggedright A subtree state is independent of environment functions, simplifying interpretation. \tabularnewline
\raggedright Exact TTN contractions
& \raggedright Need joint mixed moments, Monte Carlo, or a tractable model of the data law.
& \raggedright Second-moment message passing and exact leafwise contractions. \tabularnewline
\raggedright Tied copies $x^{\otimes N}$
& \raggedright Exact neural-network equivalence for structured tied cores; use the moment operator, symmetric/Hermite output modes, and computational HSVD as a computation diagnostic.
& \raggedright Not an independent-block specialization: the formal leaves remain identical. \tabularnewline
\hline
\end{tabular}
\caption{The independence assumption is a source of orthogonality and computational closure, not a prerequisite for comparing TTN and DNN training dynamics.}
\label{tab:evidence}
\end{table}

The main methodological conclusion is therefore two-tiered.  Use independent blocks first when the
goal is to validate the measurement apparatus and obtain a clean, additive decomposition of learning.
Then keep the same probes, output projections and interventions while moving to dependent physical
inputs and to the tied repeated-input endpoint.  At that endpoint, success should mean agreement in
polynomial output modes, hidden/computational feature subspaces and intervention responses--not an
unqualified identification of copy-cut singular values with canonical features of the input--output
function.
